\section*{Capítulo \ref{ch:b1:extent-q-and-k} --- Descrever um sistema químico: avanço, atividades, Q e K}

\begin{solution}{exo:b1:extent-q-and-k:1}
\omterm{def:b1:extent-q-and-k:variables}{Intensivas}: temperatura, \omterm{def:b1:extent-q-and-k:composition}{concentração molar}, massa específica, pressão,
\omterm{def:b1:extent-q-and-k:composition}{fração molar}. \omterm{def:b1:extent-q-and-k:variables}{Extensivas}: volume, quantidade de matéria, massa.
\end{solution}

\begin{solution}{exo:b1:extent-q-and-k:2}
$p(\ce{N2}) = 0.7808 \times 1.013 = \qty{0.791}{bar}$, $p(\ce{O2}) =
0.2095 \times 1.013 = \qty{0.212}{bar}$, $p(\ce{Ar}) = 0.0093 \times 1.013
= \qty{0.0094}{bar}$.
\end{solution}

\begin{solution}{exo:b1:extent-q-and-k:3}
$n(\ce{H2}) = 3.0 - 2\xi$, $n(\ce{O2}) = 2.0 - \xi$, $n(\ce{H2O}) =
2\xi$. $\xi_{\max} = \min(3.0/2,\ 2.0/1) = \qty{1.5}{mol}$: o di-hidrogênio
é o limitante. Estado final: \ce{H2} 0, \ce{O2} \qty{0.5}{mol}, \ce{H2O}
\qty{3.0}{mol}.
\end{solution}

\begin{solution}{exo:b1:extent-q-and-k:4}
(a) $Q = p_{\ce{CO2}}/p^\circ$. (b) $Q = (c^\circ)^2/([\ce{Ag+}][\ce{Cl-}])$.
(c) $Q = (p_{\ce{NH3}}/p^\circ)^2/[(p_{\ce{N2}}/p^\circ)(p_{\ce{H2}}/p^\circ)^3]$.
(d) $Q = [\ce{NH4+}][\ce{OH-}]/([\ce{NH3}]\,c^\circ)$, pois a água tem \omterm{def:b1:extent-q-and-k:activity}{atividade} 1.
\end{solution}

\begin{solution}{exo:b1:extent-q-and-k:5}
$\mathrm{A \rightleftharpoons C}$ é a soma: $K = 10 \times 0.5 = 5$. $\mathrm{C \rightleftharpoons A}$: $1/5 =
0.2$. $\mathrm{2\,A \rightleftharpoons 2\,C}$: $5^2 = 25$.
\end{solution}

\begin{solution}{exo:b1:extent-q-and-k:6}
$Q = 5.0 \times 10^{-5} < K^\circ$: sentido direto. $Q = 0.10 > K^\circ$:
sentido inverso. Só reagentes: $Q = 0 < K^\circ$, sentido direto.
\end{solution}

\begin{solution}{exo:b1:extent-q-and-k:7}
$n(\ce{A2}) = 1 - \xi$, $n(\ce{A}) = 2\xi$, total $1 + \xi$. Com $p =
p^\circ$, $Q = \dfrac{(2\xi/(1+\xi))^2}{(1-\xi)/(1+\xi)} = \dfrac{4\xi^2}{1 -
\xi^2}$. $Q = 0.25$ dá $4.25\,\xi^2 = 0.25$, $\xi = \qty{0.243}{mol}$.
\end{solution}

\begin{solution}{exo:b1:extent-q-and-k:8}
Com $x = [\ce{C}]$: $x/(0.10 - x)^2 = 100$, logo $100x^2 - 21x + 1 = 0$ e
$x = (21 - \sqrt{41})/200 = \qty{0.0730}{mol/L}$ (a outra raiz é maior que
0.10). $[\ce{A}] = [\ce{B}] = \qty{0.0270}{mol/L}$, $[\ce{C}] =
\qty{0.0730}{mol/L}$, $\tau = 0.73$.
\end{solution}

\begin{solution}{exo:b1:extent-q-and-k:9}
A partir de $n$ de cada um: $Q = \xi^2/(n - \xi)^2 = K^\circ$, logo $\xi/(n - \xi) =
\sqrt{K^\circ}$ e $\tau = \xi/n = \sqrt{K^\circ}/(1 + \sqrt{K^\circ})$.
$\tau = 0.99$ exige $\sqrt{K^\circ} = 99$, $K^\circ \approx \num{9.8e3}$:
cerca de $10^4$, o limiar usual de uma \omterm{def:b1:extent-q-and-k:quantitative}{reação quantitativa}.
\end{solution}

\begin{solution}{exo:b1:extent-q-and-k:10}
No equilíbrio, $p_{\ce{CO2}} = 0.20\,p^\circ$, ou seja, $n(\ce{CO2}) =
pV/RT = 0.20 \times 10^5 \times 0.0100/(8.314 \times 1100) =
\qty{0.0219}{mol}$. Com \qty{0.010}{mol} de carbonato, isso não pode
ser atingido: tudo se decompõe, $p_{\ce{CO2}} = 0.010 \times 8.314 \times
1100/0.0100 = \qty{9.15e3}{Pa} = \qty{0.091}{bar}$, $Q < K^\circ$, e o
estado final (\qty{0.010}{mol} de \ce{CaO}, \qty{0.010}{mol} de \ce{CO2},
nenhum \ce{CaCO3}) não é um equilíbrio. Com \qty{0.050}{mol}: equilíbrio,
\qty{0.0219}{mol} de \ce{CO2} e de \ce{CaO}, restando \qty{0.0281}{mol} de \ce{CaCO3}.
\end{solution}

\begin{solution}{exo:b1:extent-q-and-k:11}
No equilíbrio, $n(\ce{B}) = 2\,n(\ce{A})$ e $n(\ce{C}) = 3\,n(\ce{A})$
(o volume se cancela). Então $6\,n(\ce{A}) = 1.0$: $n(\ce{A}) =
\qty{0.167}{mol}$, $n(\ce{B}) = \qty{0.333}{mol}$, $n(\ce{C}) =
\qty{0.500}{mol}$.
\end{solution}

\begin{solution}{exo:b1:extent-q-and-k:12}
Quantidades iguais: $\xi^2/(0.50 - \xi)^2 = 4$, $\xi/(0.50 - \xi) = 2$,
$\xi = \qty{0.333}{mol}$, \omterm{def:b1:extent-q-and-k:fractional-extent}{rendimento} $0.333/0.50 = 67\,\%$. Com \qty{1.0}{mol}
de álcool: $\xi^2 = 4(0.50 - \xi)(1.0 - \xi)$, $3\xi^2 - 6\xi + 2 = 0$,
$\xi = (6 - \sqrt{12})/6 = \qty{0.423}{mol}$, \omterm{def:b1:extent-q-and-k:fractional-extent}{rendimento} de 85\,\%: um
excesso de um reagente aumenta o \omterm{def:b1:extent-q-and-k:fractional-extent}{rendimento} no produto do outro.
\end{solution}

\begin{solution}{pb:b1:extent-q-and-k:1}
\textbf{1.} $x(\ce{CH4}) = 1.00/4.00 = 0.250$; $x(\ce{H2O}) = 0.750$.
\textbf{2.} \qty{7.5}{bar} e \qty{22.5}{bar}.
\textbf{3.} Massas de 16.04 e \qty{54.06}{g}: $w(\ce{CH4}) = 0.229$,
$w(\ce{H2O}) = 0.771$.
\textbf{4.} \omterm{def:b1:extent-q-and-k:composition}{Frações molares}, \omterm{def:b1:extent-q-and-k:composition}{frações mássicas}, \omterm{def:b1:extent-q-and-k:composition}{pressões parciais} (e a
pressão total); as quantidades de matéria são \omterm{def:b1:extent-q-and-k:variables}{extensivas}.
\textbf{5.} $a = p_i/p^\circ$: 7.5 para o metano, 22.5 para o vapor.
\textbf{6.} \ce{CH4} $1.00 - \xi$, \ce{H2O} $3.00 - \xi$, \ce{CO} $\xi$,
\ce{H2} $3\xi$, total $4.00 + 2\xi$.
\textbf{7.} $\xi_{\max} = \qty{1.00}{mol}$; o metano é o limitante.
\textbf{8.} \ce{CH4} 0, \ce{H2O} \qty{2.00}{mol}, \ce{CO} \qty{1.00}{mol},
\ce{H2} \qty{3.00}{mol}.
\textbf{9.} \qty{6.00}{mol}, contra \qty{4.00}{mol} na alimentação: a
reação forma quatro moléculas a partir de duas.
\textbf{10.} \ce{H2O} 0.333, \ce{CO} 0.167, \ce{H2} 0.500.
\textbf{11.} Para garantir que todo o metano reaja e deixar vapor para a
reação de deslocamento, cujo quociente ele diminui.
\textbf{12.} \ce{CO} $1 - \xi$, \ce{H2O} $2 - \xi$, \ce{CO2} $\xi$, \ce{H2}
$3 + \xi$, total 6.00 em qualquer avanço.
\textbf{13.} $Q = \dfrac{(p_{\ce{CO2}}/p^\circ)(p_{\ce{H2}}/p^\circ)}
{(p_{\ce{CO}}/p^\circ)(p_{\ce{H2O}}/p^\circ)}$ com $p_i = (n_i/6)p$: os
fatores $p/(6p^\circ)$ se cancelam entre numerador e denominador (duas
moléculas de gás de cada lado), restando a expressão dada.
\textbf{14.} $\xi = 0$: $Q = 0 < 4.0$, sentido direto.
\textbf{15.} $\xi(3 + \xi) = 4(1 - \xi)(2 - \xi)$, ou seja, $3\xi^2 - 15\xi +
8 = 0$: $\xi = (15 - \sqrt{129})/6 = \qty{0.607}{mol}$; a outra raiz,
4.39, é maior que $\xi_{\max} = 1$.
\textbf{16.} \ce{CO} \qty{0.393}{mol}, \ce{H2O} \qty{1.393}{mol},
\ce{CO2} \qty{0.607}{mol}, \ce{H2} \qty{3.607}{mol}.
\textbf{17.} $\tau = 0.607/1.00 = 0.607$.
\textbf{18.} $\xi(3 + \xi) = 4(1 - \xi)(5 - \xi)$, $3\xi^2 - 27\xi + 20 =
0$, $\xi = \qty{0.814}{mol}$: mais vapor, mais conversão.
\textbf{19.} $K^\circ = (0.99 \times 3.99)/(0.01 \times 1.01) = 391$.
\textbf{20.} $Q(\xi)$ é crescente, de modo que um $K^\circ$ maior é
atingido com um avanço maior. A constante dessa reação é maior a
temperatura mais baixa, onde a reação é lenta: um primeiro reator
trabalha quente e rápido, um segundo, mais frio, completa a conversão (a
dependência de $K^\circ$ com a temperatura é tratada no volume do 2.º ano).
\textbf{21.} $3.607/(6.00 - 1.393) = 3.607/4.607 = 0.783$.
\textbf{22.} \qty{0.393}{mol} de \ce{CO} e \qty{0.607}{mol} de \ce{CO2}.
\textbf{23.} $x(\ce{H2}) = 3.607/6.00 = \textbf{0.60}$: seis moléculas em
cada dez que saem do reator de deslocamento são de hidrogênio.
\end{solution}
